\chapter{Cryptography — Arithmetic Circuit}
%%% generated by the blueprint pipeline from TCSlib.Cryptography.MPC.ArithmeticCircuit
\section{Overview}

Declarations of the Lean module \texttt{TCSlib.Cryptography.MPC.ArithmeticCircuit}.

\section{Declarations}

\begin{definition}[Arithmetic circuit syntax]
\lean{TCSlib.MPC.Circuit}
\statementsource{BGW88}{
  pdf-b3f55f2a4332-p003-b004,
  pdf-b3f55f2a4332-p003-b008
}
\statementsource{AL17}{
  pdf-27f2b3787521-p003-b007,
  pdf-27f2b3787521-p015-b006,
  pdf-27f2b3787521-p018-b003
}
\label{TCSlib.MPC.Circuit}
\leanok
For a coefficient type $F$ and input-index type $\iota$, an arithmetic circuit is
generated by input wires indexed by $\iota$, public constants from $F$, sums of two
circuits, scalar multiples of a circuit by a public coefficient, and products of two
circuits.
\end{definition}

\begin{definition}[Evaluation of an arithmetic circuit]
\lean{TCSlib.MPC.Circuit.eval}
\statementsource{BGW88}{
  pdf-b3f55f2a4332-p003-b004,
  pdf-b3f55f2a4332-p003-b008
}
\statementsource{AL17}{
  pdf-27f2b3787521-p003-b007,
  pdf-27f2b3787521-p018-b003
}
\label{TCSlib.MPC.Circuit.eval}
\leanok
\uses{TCSlib.MPC.Circuit}
Let $F$ be a commutative ring, let $\iota$ be an input-index type, and assign a value in
$F$ to each input index. The value of a circuit is obtained by interpreting input wires
by this assignment, constants by their given values, addition and multiplication gates
by the ring operations, and scalar multiplication by multiplication by its public
coefficient.
\end{definition}

\begin{theorem}[Evaluation of an Input Gate]
\lean{TCSlib.MPC.Circuit.eval_input}
\label{TCSlib.MPC.Circuit.eval_input}
\leanok
\uses{TCSlib.MPC.Circuit, TCSlib.MPC.Circuit.eval}
\difficulty{1}
Let $F$ be a commutative ring, let $v$ assign an element of $F$ to each input index, and
let $i$ be an input index. The circuit consisting of the input wire $i$ evaluates to
$v(i)$.
\end{theorem}

\begin{theorem}[Evaluation of a Constant Gate]
\lean{TCSlib.MPC.Circuit.eval_const}
\label{TCSlib.MPC.Circuit.eval_const}
\leanok
\uses{TCSlib.MPC.Circuit, TCSlib.MPC.Circuit.eval}
\difficulty{1}
Let $F$ be a commutative ring, let $v$ assign elements of $F$ to the input indices, and
let $c\in F$. The circuit consisting of the constant $c$ evaluates to $c$.
\end{theorem}

\begin{theorem}[Evaluation of an Addition Gate]
\lean{TCSlib.MPC.Circuit.eval_add}
\label{TCSlib.MPC.Circuit.eval_add}
\leanok
\uses{TCSlib.MPC.Circuit, TCSlib.MPC.Circuit.eval}
\difficulty{1}
Let $F$ be a commutative ring, let $v$ assign an element of $F$ to each input index, and
let $a,b$ be arithmetic circuits. Evaluating the circuit formed by adding $a$ and $b$
gives the sum of their evaluations: $(a+b)(v)=a(v)+b(v)$.
\end{theorem}

\begin{theorem}[Evaluation of scalar multiplication]
\lean{TCSlib.MPC.Circuit.eval_smul}
\proofsource{AL17}{pdf-27f2b3787521-p016-b002}
\label{TCSlib.MPC.Circuit.eval_smul}
\leanok
\proofstep
  {TCSlib.MPC.Circuit}
  {context}
  {AL17}
  {pdf-27f2b3787521-p015-b006,
    pdf-27f2b3787521-p016-b001}
\proofstep
  {TCSlib.MPC.Circuit.eval}
  {context}
  {AL17}
  {pdf-27f2b3787521-p018-b003}
\proofstep
  {TCSlib.MPC.Circuit.eval_input}
  {context}
  {AL17}
  {pdf-27f2b3787521-p018-b003}
\proofstep
  {TCSlib.MPC.Circuit.eval_const}
  {context}
  {AL17}
  {pdf-27f2b3787521-p015-b006}
\proofstep
  {TCSlib.MPC.Circuit.eval_add}
  {context}
  {AL17}
  {pdf-27f2b3787521-p016-b002}
\proofstep
  {TCSlib.MPC.Circuit.eval_mul}
  {context}
  {AL17}
  {pdf-27f2b3787521-p016-b003,
    pdf-27f2b3787521-p018-b003}
\uses{TCSlib.MPC.Circuit, TCSlib.MPC.Circuit.eval}
\difficulty{1}
Let $F$ be a commutative ring, let $\iota$ be a type of input wires, and assign each
input $i$ a value $v(i)\in F$. For every circuit $a$ and scalar $c\in F$, the circuit
obtained by multiplying $a$ by the public constant $c$ evaluates to $c$ times the value
of $a$.
\end{theorem}

\begin{theorem}[Evaluation of a Multiplication Gate]
\lean{TCSlib.MPC.Circuit.eval_mul}
\label{TCSlib.MPC.Circuit.eval_mul}
\leanok
\uses{TCSlib.MPC.Circuit, TCSlib.MPC.Circuit.eval}
\difficulty{1}
Let $F$ be a commutative ring, let $v$ assign an element of $F$ to each input index, and
let $a,b$ be arithmetic circuits. Evaluating the circuit formed by multiplying $a$ and
$b$ gives the product of their evaluations: $(ab)(v)=a(v)b(v)$.
\end{theorem}

\begin{definition}[Syntactically linear circuits]
\lean{TCSlib.MPC.Circuit.IsLinear}
\statementsource{AL17}{
  pdf-27f2b3787521-p021-b004,
  pdf-27f2b3787521-p022-b001
}
\label{TCSlib.MPC.Circuit.IsLinear}
\leanok
\uses{TCSlib.MPC.Circuit}
Let $F$ be a commutative ring and let $\iota$ be an input-index type. Input wires and
constants are linear circuits; a sum is linear exactly when both of its inputs are
linear; scalar multiplication preserves linearity; and every multiplication gate is
classified as nonlinear.
\end{definition}
